\section{Experimental protocol}
\label{sec:setup}

\subsection{Splits}
\label{sec:splits}

\paragraph{Subject-wise folds}
Subjects are shuffled once and partitioned into five blocks. Fold $f$ tests
on block $f$, calibrates on four subjects drawn from the remaining blocks
(rotated across folds), and trains on the rest. The test blocks partition the
cohort, so every subject contributes exactly one test set across the five
folds, and no test epoch shares a subject with any calibration or training
epoch of the same fold. Per-subject statistics are computed by merging the
five folds and then averaging each subject's coverage over the seeds and
backbones that re-run that partition. Each subject therefore contributes
exactly one coverage value to the worst-patient and dispersion statistics,
and the spread we report for them is the run-to-run spread of the training,
not a spread over overlapping fold summaries.

The four-subject calibration cohort is deliberate: it is the regime a
hospital faces when it has a handful of scored nights, and it is where the
marginal guarantee is weakest. \Cref{sec:results-cohort-size} varies the size
explicitly.

\paragraph{Random-epoch split}
As a diagnostic we also train and calibrate with a 60/20/20 epoch-level
random split. Calibration and test epochs then come from the same subjects,
which is the protocol under which most published epoch-level numbers are
computed. It is not a valid estimate of deployment performance, and we use it
only to show how much of the coverage gap it hides.

\paragraph{Cross-cohort}
For the shift experiment, models are trained and calibrated on SC subjects
and tested on ST subjects, and vice versa, with no subject overlap in either
direction.

\subsection{Models and training}
\label{sec:models}

Two backbones are used throughout so that conclusions are not tied to a
single architecture: \texttt{eegnet}, a depthwise-separable convolutional
network in the style of \citet{lawhern2018} with \num{1445} parameters, and
\texttt{sleepcnn}, a four-block one-dimensional VGG-style network with global
average pooling and \num{224069} parameters. Both take a
$3000$-sample epoch and output a five-way softmax.

Training uses AdamW (learning rate $10^{-3}$, weight decay $10^{-4}$), a
cosine schedule, batch size 128, at most 40 epochs, and early stopping on a
held-out subset of the training subjects. The loss is plain cross-entropy
without class weights in every run reported here. Class weighting trades N1
recall against the calibration of the output distribution, which is exactly
the quantity CP consumes, and mixing the two settings across tables would
make coverage numbers incomparable; we therefore fix one setting and record
it in each run's metadata.

In the main matrix every configuration is run with three seeds (0, 1, 2) and
five folds, giving $15$ runs per backbone and $30$ in total; the extended
cohort and the cross-cohort experiment use seed 0 only, so their numbers are
single-seed estimates over five folds and two backbones, and the spread
reported there is across those ten runs rather than across seeds. The
run-to-run spread over seeds and backbones is
reported as the uncertainty on every headline number of the main matrix, and
per-subject statistics average a subject's coverage over those runs before
any worst-patient or dispersion statistic is formed.

\paragraph{Hardware and software}
All experiments run on a single NVIDIA RTX 4060 Laptop GPU with PyTorch;
preprocessing, conformal procedures and dynamic programs are NumPy and run on
one CPU core. Training is seeded and cuDNN is pinned to deterministic kernels,
so repeating a training run reproduces its test probabilities exactly; we
verified this by retraining one fold while the rest of the matrix occupied the
GPU and comparing the arrays. The released code reproduces every table and
figure from the raw archives with the commands recorded in the repository
README.

\subsection{Conformal procedures and hyperparameters}
\label{sec:conformal-setup}

Unless stated otherwise $\alpha=0.1$. All thresholds are computed on
calibration subjects only. The procedures compared are:

\begin{itemize}
\item \emph{Pooled APS} and \emph{pooled RAPS} on the raw softmax, with the
      split-conformal threshold \Cref{eq:split-cp}.
\item \emph{Smoothed APS}: APS on the forward--backward posterior
      \Cref{eq:forward-backward}.
\item \emph{Smoothed deficit}: the score \Cref{eq:deficit} on the same
      posterior.
\item \emph{Per-stage thresholds}: the Mondrian construction of
      \Cref{sec:scores}, applied to both the raw APS score and the smoothed
      deficit score, so that the stage-level comparison uses the same scores
      as the pooled one.
\item \emph{SACP}: warm-up calibration with shrinkage \Cref{eq:shrunk-q}, for
      label budgets $m\in\{0,10,20,40,80,160,320\}$ and $k\in\{3,5,10\}$
      neighbours for the label-free variant, with $\lambda$ chosen by
      leave-one-subject-out simulation on the calibration subjects.
\item \emph{Per-patient ACI}: one stream per patient \Cref{eq:aci-patient},
      with $\gamma_0$ chosen by leave-one-subject-out simulation from
      $\{0.005,0.01,0.02,0.05,0.1\}$.
\item \emph{SeqCP-Margin} and \emph{SeqCP-Block}: the sequence constructors of
      \Cref{sec:method-seq}, for block lengths $L\in\{10,20,30\}$ and
      transition weights $\beta\in\{0,0.5,1,2\}$.
\item \emph{Oracle}: per-subject thresholds computed on the test subject's
      own labels. Not deployable; it is the lower bound on set size at
      nominal coverage and shows how much room per-patient adaptation has.
\end{itemize}

HMM transitions are estimated from the training labels with Laplace
smoothing. Block-level comparisons restrict all methods to exactly the epochs
inside complete blocks, and the per-epoch baselines in those comparisons are
recalibrated on the same calibration epochs, so the block-length comparison
isolates the construction rather than the calibration sample.

\subsection{Metrics}
\label{sec:metrics}

Coverage is reported marginal, per stage, and per subject; from the
per-subject values we report the mean, the 10th percentile, the worst subject
and the fraction of subjects below nominal. Efficiency is reported as mean set
size, singleton rate and review workload, the fraction of epochs whose set
contains more than one stage and would therefore need adjudication. Temporal
behaviour is reported with the
Jaccard consistency statistic of \Cref{sec:notation} and with sequence
coverage, the fraction of complete blocks all of whose epochs are covered.
Classification quality is reported with accuracy, Cohen's kappa, macro-F1 and
expected calibration error (15 equal-mass bins), so that the coverage results
can be read against the underlying model quality.

